\documentclass[11pt]{article}
\usepackage[margin=1in]{geometry}
\usepackage{mathpazo}
\usepackage{amsmath,amssymb}
\usepackage{microtype}
\usepackage{booktabs}
\usepackage{graphicx}
\usepackage[dvipsnames]{xcolor}
\usepackage{tikz}
\usetikzlibrary{positioning,arrows.meta,calc,fit,backgrounds}
\usepackage[most]{tcolorbox}
\usepackage{enumitem}
\usepackage{listings}
\usepackage[font=small,labelfont=bf]{caption}
\usepackage{hyperref}
\hypersetup{colorlinks=true,linkcolor=MidnightBlue,citecolor=MidnightBlue,urlcolor=MidnightBlue}
\usepackage{cleveref}
\urlstyle{tt}
\def\UrlBreaks{\do\_\do/\do.}
\newcommand{\file}[1]{\nolinkurl{#1}}
\newcommand{\code}[1]{\texttt{#1}}
\newcommand{\mhh}{m_{hh}}
\newcommand{\bw}{\boldsymbol{\omega}}
\newcommand{\bp}{\boldsymbol{p}}
\newcommand{\bx}{\boldsymbol{x}}
\newcommand{\data}{\mathcal{D}}

\definecolor{trainable}{HTML}{1B6CA8}
\definecolor{physics}{HTML}{2E933C}
\definecolor{datac}{HTML}{7B4EA3}
\definecolor{warn}{HTML}{C0392B}

\newtcolorbox{keyidea}[1][]{enhanced,breakable,colback=physics!6,colframe=physics!70!black,
  fonttitle=\bfseries,title={#1},boxrule=0.6pt,arc=2pt,left=6pt,right=6pt}
\newtcolorbox{aside}[1][]{enhanced,breakable,colback=trainable!5,colframe=trainable!70!black,
  fonttitle=\bfseries,title={#1},boxrule=0.6pt,arc=2pt,left=6pt,right=6pt}
\newtcolorbox{assume}[1][]{enhanced,breakable,colback=datac!5,colframe=datac!70!black,
  fonttitle=\bfseries,title={#1},boxrule=0.6pt,arc=2pt,left=6pt,right=6pt}
\newtcolorbox{watch}[1][]{enhanced,breakable,colback=warn!5,colframe=warn!70!black,
  fonttitle=\bfseries,title={#1},boxrule=0.6pt,arc=2pt,left=6pt,right=6pt}
\setlist{itemsep=2pt,topsep=3pt}
\setlength{\emergencystretch}{3em}

\lstset{language=Python, basicstyle=\ttfamily\small, columns=fullflexible,
  keepspaces=true, commentstyle=\color{physics!70!black}\itshape,
  keywordstyle=\color{trainable!90!black}\bfseries, stringstyle=\color{datac},
  showstringspaces=false, frame=single, rulecolor=\color{gray!40},
  backgroundcolor=\color{gray!3}, xleftmargin=4pt, xrightmargin=4pt,
  aboveskip=6pt, belowskip=6pt}

\title{\textbf{A bootstrap ensemble of \texttt{HEFTNet} models}\\[4pt]
\large Plan for \texttt{heft\_train\_ensemble.ipynb} and \texttt{dihiggs/ensemble.py}\\[4pt]
\normalsize Revised plan: bootstrap of the training rows. For approval before any code is written}
\date{October 30, 2026 (supersedes the first version of the same date)}

\begin{document}
\maketitle

\begin{abstract}
\noindent
The first Bayesian neural network run (\file{docs/bnn/bnn_run1.pdf}) gave credible
bands about 17 times narrower than the model's actual error, partly because a single
HMC chain explores a single valley of the posterior. This plan describes a
complementary, frequentist approach: train many copies of \code{HEFTNet}, each on a
\emph{bootstrap} resample of the training rows and with its own random seed, and use the
spread of their predictions as an estimate of the standard error of the trained network.
A second ensemble, which varies only the seed, separates the part of that spread due to
the optimisation (which valley) from the part due to the finite training sample. The
first version of this plan resampled whole spectra; following our discussion, it now
resamples the training data themselves, as the pipeline sees them. Which resampling is
``right'' is an instance of a classic and unresolved problem of statistical inference,
the choice of the \emph{reference ensemble}, discussed in \cref{sec:reference}. The work
is a new notebook, \file{code/heft_train_ensemble.ipynb}, and a small new module,
\file{dihiggs/ensemble.py}. Nothing will be written until you approve; the decisions
needed are collected in \cref{sec:decisions}.
\end{abstract}

\begin{keyidea}[The proposal in brief]
\begin{enumerate}
  \item \textbf{Two ensembles} of $M = 20$ members each (\cref{tab:ensembles}):
    \code{'seeds'} (the training rows as they are, different seeds) and
    \code{'bootstrap'} (the training rows resampled with replacement, different seeds).
  \item \textbf{The bootstrap embeds our sample in the population it was drawn from,}
    as seen by the pipeline: the training rows are regarded as independent draws from a
    (discrete, complicated) distribution $P$, and each bootstrap set is a draw of the
    same size from the empirical distribution of those rows.
  \item \textbf{The split of \file{heft_train.ipynb} is kept exactly:} the same training,
    validation and quick-test rows. Only the training rows are resampled; the
    validation rows are the same for every member.
  \item \textbf{Each member is trained exactly as in \file{heft_train.ipynb}} and saved as an
    ordinary run folder, so \file{heft_test.ipynb} and \file{heft_export.ipynb} work with
    any member.
  \item \textbf{An \code{Ensemble} class} with the same methods as \code{hmc.BNN}.
  \item \textbf{The test,} on the 1,000 holdout spectra: is the ensemble spread, added to
    the POWHEG noise, the size of the actual differences from POWHEG?
  \item \textbf{The reference-ensemble problem} (\cref{sec:reference}): the answer depends
    on the population into which the sample is embedded, and there is no unique choice.
    The plan states its choice explicitly.
\end{enumerate}
\end{keyidea}

\tableofcontents

% ===========================================================================
\section{Why an ensemble}
\label{sec:why}

The BNN run showed three reasons why a single chain's band is too narrow
(\file{bnn_run1.pdf}, \S6): incomplete exploration of the valley it is in; other valleys,
in which equally good networks give different predictions where the data constrain them
weakly; and the model's own systematic error, which no spread of parameters within one
model describes. An ensemble of independently trained networks addresses the second
directly. Each training run, started from different random weights and fed batches in a
different order, ends in a different valley. Large studies have found that such
``deep ensembles'' reproduce the predictions of long, full-data HMC better than most
other approximations, precisely because they visit several valleys~\cite{izmailov2021,
wilson2020, lakshminarayanan2017}. Resampling the training data as well, by the
bootstrap~\cite{efron1979, efron1993}, adds the variation that any estimator has because
its training set is one finite sample.

% ===========================================================================
\section{What the ensemble estimates}
\label{sec:what}

\subsection{The frequentist view}

Think of the whole training procedure as an \emph{estimator}: it takes a training set
$D$ and a random seed $s$ and returns a function $\hat f(\bx; D, s)$. Had the training set
been a different draw from the same population, or had the seed been different, the
result would have been different. So $\hat f(\bx)$ at a fixed point $\bx$ is a random
variable, with a mean and a variance. Its standard deviation is the \emph{standard error}
of the trained network at $\bx$, the frequentist counterpart of the width of a Bayesian
band.

\begin{aside}[The bootstrap principle]
We have only one training set, so the distribution of $\hat f$ over training sets cannot
be observed directly. The bootstrap~\cite{efron1979, efron1993} replaces the unknown
distribution $P$ from which the $N$ training rows were drawn by their \emph{empirical}
distribution $\hat P$, which puts probability $1/N$ on each row. A bootstrap training set
$D^*$ is a sample of $N$ rows drawn from $\hat P$, i.e.\ drawn from the training rows
\emph{with replacement}. On average it contains $1 - (1 - 1/N)^N \approx 1 - e^{-1}
\approx 63.2\%$ of the distinct rows, some of them several times. The spread of
$\hat f(\bx; D^*_m, s_m)$ over $m = 1, \ldots, M$ estimates the spread of $\hat f$ over
training sets drawn from $P$.
\end{aside}

\begin{aside}[Your thought experiment]
Suppose the colleague had generated a hundred times more data, and the pipeline were run
on each of 100 disjoint pieces of it. The spread of the 100 trained networks is the
quantity of interest: the sampling variability of the procedure at the size of data we
actually have. Each piece is a draw from $P$, and the bootstrap mimics such draws by
drawing from $\hat P$. The quality of the mimicry rests on $\hat P$ being close to $P$,
which, with 274,152 training rows, is a reasonable expectation for the smooth functionals
that a trained network represents.
\end{aside}

\subsection{Separating the two sources}
\label{sec:decomp}

The law of total variance splits the variance of $\hat f$ at $\bx$ into an optimisation
part and a data part:
\begin{equation}
  \underbrace{\operatorname{Var}_{D,s}\bigl[\hat f\bigr]}_{\text{total}}
  \;=\;
  \underbrace{\operatorname{E}_{D}\Bigl[\operatorname{Var}_{s}\bigl[\hat f \mid D\bigr]\Bigr]}_{\text{optimisation: which valley}}
  \;+\;
  \underbrace{\operatorname{Var}_{D}\Bigl[\operatorname{E}_{s}\bigl[\hat f \mid D\bigr]\Bigr]}_{\text{data: which sample}} .
  \label{eq:totvar}
\end{equation}
The two ensembles of \cref{tab:ensembles} estimate the first term and the total. The data
term is their difference, assuming that the optimisation spread is about the same for a
bootstrap set as for the original set.

\begin{table}[h]
\centering\small
\begin{tabular}{lllll}
\toprule
\code{MODE} & training rows of member $m$ & seed & estimates & run folder\\
\midrule
\code{'seeds'} & the 274,152 training rows & \code{SEED} $+ m$ & optimisation term & \file{runs/ENS_seeds}\\
\code{'bootstrap'} & 274,152 rows drawn from them & \code{SEED} $+ m$ & total & \file{runs/ENS_bootstrap}\\
\bottomrule
\end{tabular}
\caption{The two ensembles. Member $m$ uses the same seed in both, so the two members
start from the same initial weights; only the training rows differ.}
\label{tab:ensembles}
\end{table}

\begin{watch}[What neither ensemble captures: the shared bias]
All members have the same architecture, the same loss and the same training schedule. If
they all miss the plateau just above threshold in the same way, they agree with each
other, and the spread stays small however many members there are. The bootstrap standard
error is an honest error bar on the \emph{estimator} $\hat f$, not on the true cross
section, unless the bias $\operatorname{E}[\hat f] - \sigma_\text{true}$ is small compared
with it. Near threshold it may not be. The test of \cref{sec:test} is designed to show
whether it is.
\end{watch}

\subsection{How many members?}

The standard deviation of $M$ draws is itself uncertain, by a relative amount of about
$1/\sqrt{2(M - 1)}$:

\begin{center}
\small
\begin{tabular}{lcccc}
\toprule
$M$ & 10 & 20 & 50 & 100\\
\midrule
relative uncertainty of the standard deviation & 24\% & 16\% & 10\% & 7\%\\
\bottomrule
\end{tabular}
\end{center}

\noindent I propose $M = 20$ per ensemble, enough to tell a factor of two in width, which
is the level of the questions we are asking. Members can be added later without redoing
the others.

% ===========================================================================
\section{The reference-ensemble problem}
\label{sec:reference}

\subsection{The problem}

Every frequentist statement, a standard error, a confidence level, a $p$-value, is a
statement about an \emph{ensemble}: the imagined infinite collection of repetitions of
the experiment of which the one we performed is regarded as a member. The data do not
determine that ensemble. A single sample can be embedded in many populations, each
consistent with how it was obtained, and each giving a different answer. This is the
statistical form of the \emph{reference class problem}, as old as the frequency theory
of probability itself: Venn~\cite{venn1866} noted that an individual belongs to many
classes, in which the frequency of an attribute differs, and Reichenbach~\cite{reichenbach1949}
proposed, without a general prescription for finding it, to take the narrowest class for
which reliable statistics exist. H\'ajek~\cite{hajek2007} argues that every
interpretation of probability, the Bayesian included, faces a version of the problem.

Within statistics proper, Fisher~\cite{fisher1956} insisted that the reference set must
contain no \emph{recognisable subset}: no subset, identifiable from the data, in which the
frequencies differ. Cox's famous example~\cite{cox1958} makes the point. A measurement is
made with one of two instruments, one precise and one imprecise, chosen by the toss of a
coin. The ensemble of all repetitions, in which either instrument may be used, gives a
valid confidence interval for the unconditional coverage; but once we know which
instrument was used, the relevant ensemble is surely the repetitions with \emph{that}
instrument, and the unconditional interval is too wide for one and too narrow for the
other. Buehler~\cite{buehler1959} and Robinson~\cite{robinson1975} turned Fisher's remark
into a criterion and showed that standard confidence intervals can have
\emph{relevant subsets} in which their coverage is systematically different. The
conditionality principle, made formal by Birnbaum~\cite{birnbaum1962} (see also
\cite{berger1988}), says one should condition on such ancillary information; carried to
its conclusion, together with sufficiency, it leads to the likelihood principle, under
which the ensemble disappears altogether, which is precisely what the frequentist cannot
accept. The problem is therefore not a technical gap but a fault line between schools of
inference; Cox and Hinkley~\cite{coxhinkley1974} and, for physicists,
Cousins~\cite{cousins1995} give balanced accounts.

\begin{aside}[An analogy from statistical mechanics]
The same system can be described by the microcanonical, canonical or grand canonical
ensemble, depending on what is held fixed: the energy, the temperature, the particle
number. The ensembles agree on some quantities and not on others, notably fluctuations.
The choice is not dictated by the system but by the question asked and by what is
regarded as fixed in the imagined repetitions. Reference ensembles in statistics behave in
the same way.
\end{aside}

\subsection{How it shows up in the bootstrap}

The bootstrap does not escape the problem; it makes it concrete. The resampling scheme
\emph{is} the choice of reference ensemble. For regression, the classic distinction is
between resampling the \emph{pairs} $(\bx_i, t_i)$, which treats the inputs as random,
and resampling the \emph{residuals} at fixed inputs, which treats them as
fixed~\cite{freedman1981, wu1986, davison1997}. Freedman showed that both are valid,
under their own assumptions, and that they estimate different variances when the model is
wrong or the noise heteroscedastic; Davison and Hinkley's guiding rule is that the
resampling scheme should mimic the way the data were generated~\cite{davison1997},
which presupposes an answer to the question of which features of that way are to be held
fixed. For our data at least three embeddings are defensible:

\begin{center}
\small
\begin{tabular}{p{0.21\textwidth}p{0.33\textwidth}p{0.38\textwidth}}
\toprule
resample & imagined repetition & regarded as fixed\\
\midrule
\textbf{training rows} (this plan) & new independent draws of rows $(\bx, t)$ from $P$ &
  the population $P$ of rows, as the pipeline sees them\\
whole spectra & a new set of HEFT points $\theta$, each with a new POWHEG run &
  the distribution from which $\theta$ is drawn\\
POWHEG noise only & new events at the \emph{same} HEFT points (a residual or parametric
  bootstrap, using the known event weights $w_j$) & the HEFT points and the bins\\
\bottomrule
\end{tabular}
\end{center}

\noindent Each answers a well-posed question, and the questions differ in what is held
fixed. None is uniquely correct, and the data cannot choose between them. The row
bootstrap chosen here treats the training data as the pipeline does: a large collection of
rows from one distribution, with no further structure. It is the most direct expression
of your thought experiment, in which the pipeline is simply re-run on more data of the
same kind.

\begin{keyidea}[The working stance]
\begin{enumerate}
  \item \textbf{State the ensemble.} Every number this ensemble produces is a standard
    error with respect to repetitions in which the training rows are independent draws
    from the population $P$ of rows. The report will say so.
  \item \textbf{Calibrate empirically.} Whatever the reference ensemble, the holdout
    spectra provide a test that does not depend on it: whether the spread, combined with
    the POWHEG noise, matches the actual differences from POWHEG (\cref{sec:test}). This
    is the frequentist's own way out: coverage, checked on data not used for training.
  \item \textbf{Treat the other embeddings as sensitivity checks,} if the result invites
    them. If, for example, the bootstrap spread proves too small near threshold, the noise
    bootstrap would tell whether the POWHEG statistics alone account for more.
\end{enumerate}
\end{keyidea}

\begin{aside}[And the Bayesian?]
The Bayesian conditions on the data actually observed and needs no ensemble of
repetitions; Cox's two instruments raise no difficulty. But the problem returns in another
form: the choice of the model and the prior, and the likelihood itself, embody
assumptions about how the data arose. Our BNN, for example, treats the 274,152 rows as
conditionally independent given the parameters, which is also a statement about the
population. The posterior predictive checks of Gelman, Meng and Stern~\cite{gelman1996}
face the analogous choice explicitly: which features of the data are held fixed in the
replications.
\end{aside}

% ===========================================================================
\section{Resampling the training rows}
\label{sec:resample}

\begin{figure}[h]
\centering
\resizebox{\textwidth}{!}{%
\begin{tikzpicture}[
  font=\small, >={Stealth[length=2.2mm]},
  box/.style={draw, rounded corners=2pt, thick, align=center, inner sep=3pt,
              minimum height=11mm, text width=30mm},
  dat/.style={box, fill=datac!8, draw=datac!80!black},
  nbk/.style={box, fill=trainable!10, draw=trainable!80!black},
  res/.style={box, fill=physics!10, draw=physics!80!black},
]
\node[dat] (file) at (0, 0) {training file\\\footnotesize 317,306 rows};
\node[dat] (val)  at (4.6, 1.9) {validation + quick test\\\footnotesize 5,078 + 38,076 rows, fixed};
\node[dat] (pool) at (4.6, -1.9) {training rows\\\footnotesize 274,152};
\node[nbk] (draw) at (9.6, -1.9) {member $m$: draw 274,152 rows with replacement};
\node[res] (rows) at (14.6, -1.9) {its training set ($\approx$ 63\% distinct rows)};
\node[res] (oob)  at (14.6, 1.9) {its out-of-bag rows ($\approx$ 37\%)};
\draw[->, thick] (file) -- node[above left, font=\footnotesize]{split of \file{heft_train}} (val);
\draw[->, thick] (file) -- (pool);
\draw[->, thick] (pool) -- node[below=2pt, font=\footnotesize, align=center]{seed\\\code{BOOT\_SEED} + $m$} (draw);
\draw[->, thick] (draw) -- (rows);
\draw[->, thick, dashed] (draw) -- node[above left, font=\footnotesize]{not drawn} (oob);
\end{tikzpicture}}
\caption{How the training set of one bootstrap member is formed. The split into training,
validation and quick-test rows is that of \file{heft_train.ipynb}, the same for all
members of both ensembles.}
\label{fig:resample}
\end{figure}

\begin{itemize}
  \item \textbf{The split.} \file{heft_train.ipynb} takes the first 86.4\% of the rows of
    the (already shuffled) training file for training, the next 1.6\% for validation and
    the last 12\% as the quick test. The ensemble notebook uses exactly the same split, so
    a \code{'seeds'} member is the \file{heft_train} model with a different seed, and only
    the 274,152 training rows are resampled. The validation rows are therefore never
    copies of training rows, and the validation loss, which chooses the saved parameters,
    is unaffected by the resampling.
  \item \textbf{The draw.} Member $m$ draws 274,152 indices uniformly, with replacement,
    from the training rows, with \code{np.random.RandomState(BOOT\_SEED + m)}, whose
    sequence is frozen across numpy versions. The drawn rows are in random order, so no
    further shuffling is needed.
  \item \textbf{Bookkeeping.} Each bootstrap member's folder gets
    \file{bootstrap_counts.npz}: for each training row, how many times it was drawn
    (compressed, about 100 kB). Together with \code{BOOT\_SEED} this makes every member
    reproducible and identifies its \emph{out-of-bag} rows, those it never saw, which
    give a free extra test of that member.
  \item \textbf{Size.} A bootstrap set has as many rows as the original, so the training
    schedule needs no change.
\end{itemize}

% ===========================================================================
\section{Training the members}
\label{sec:train}

Each member is trained with the code of \file{heft_train.ipynb}, unchanged except for the
training rows: L1 loss, Adam, batch size 2,048, 200,000 iterations in four learning-rate
steps from $2 \times 10^{-3}$, parameters saved at the best validation loss. Before member
$m$ is built, \code{torch.manual\_seed(SEED + m)} and \code{np.random.seed(SEED + m)} fix
its initial weights and its batch order.

\paragraph{Run folders.} Member $m$ of an ensemble is saved in
\file{runs/ENS_<mode>/member_01/}, \file{member_02/}, \ldots, each an ordinary run folder
(\file{dihiggs_config.yaml}, \file{dihiggs_params.pth}, \file{dihiggs_losses.csv}), plus
\file{bootstrap_counts.npz} for bootstrap members. The configuration records \code{MODE},
the member number, both seeds and the training file's SHA-256. A summary file
\file{runs/ENS_<mode>/ensemble.csv} has one row per member: seeds, number of distinct
rows, best validation loss, out-of-bag loss, training time.

\paragraph{Interruption and parallel running.} Members are trained one after another, and
a member whose parameter file and summary row already exist is skipped, so an interrupted
loop resumes where it stopped. A switch \code{MEMBERS} (e.g.\ \code{range(0, 10)})
chooses which members a notebook trains, so that two or more copies of the notebook can
train disjoint sets of members at the same time on the 5090. \code{HEFTNet} is small;
several trainings fit on one GPU at once. Given the unstable connection to the 5090,
running the notebook headless (\code{jupyter nbconvert --execute} inside \code{tmux}) is
recommended.

\paragraph{The real-time loss monitor.} \file{heft_train.ipynb} opens a window that plots
the losses during training. In a loop over 20 members that would open 20 windows, so the
ensemble notebook writes the loss files without starting the window; progress is printed
instead.

\paragraph{Cost.} About 42 minutes per member on a CPU; on the 5090 I expect a few minutes.
The notebook times the first member and prints an estimate for the rest. The
\code{N\_ITERATIONS} switch allows shorter members for a first trial.

% ===========================================================================
\section{The module \texttt{dihiggs/ensemble.py}}
\label{sec:module}

A class \code{Ensemble} with the same interface as \code{hmc.BNN}, so that analysis code
can use either:

\begin{lstlisting}
from dihiggs import ensemble
from dihiggs import heftnet as NN

ens = ensemble.Ensemble(NN.HEFTNet, 'runs/ENS_bootstrap')  # every member_*/
mu  = ens.mean(x)                  # ensemble mean, one value per row of x
sd  = ens.std(x)                   # spread of the members (the standard error)
lo, hi = ens.interval(x, cl=0.68)  # central 68% of the members
A   = ens.apply(lambda m, x: m.coeffs(x), xs)   # a_i(m_hh) for every member
\end{lstlisting}

\noindent It reads \file{member_*/dihiggs_params.pth} by default (any file pattern can be
given), and also provides \code{summary()}, the contents of \file{ensemble.csv}. Like
\file{hmc.py}, it imports nothing from \code{dihiggs}, so it can move to \code{drwho}.
\code{Ensemble} is written as a general ``ensemble of parameter files''; when both modules
move to \code{drwho}, \code{hmc.BNN} can become a subclass of it. \file{hmc.py} is not
changed now.

% ===========================================================================
\section{The notebook \texttt{heft\_train\_ensemble.ipynb}}
\label{sec:notebook}

In the style of \file{heft_train.ipynb} (header, Colab cell with \code{USE\_DRIVE},
imports, device, data with its SHA-256), then:

\begin{enumerate}
  \item \textbf{Introduction} (markdown): the estimator view, the bootstrap, the
    decomposition (\ref{eq:totvar}), and a short statement of the reference ensemble
    used.
  \item \textbf{Switches:} \code{MODE = 'bootstrap'} or \code{'seeds'};
    \code{N\_MEMBERS = 20}; \code{MEMBERS = range(N\_MEMBERS)}; \code{SEED = 42};
    \code{BOOT\_SEED = 1000}; \code{N\_ITERATIONS = 200\_000}.
    \code{DIRNAME} = \file{ENS_<mode>}.
  \item \textbf{The split} of \file{heft_train.ipynb}.
  \item \textbf{One member's data:} a function returning the training rows of member $m$
    and its counts.
  \item \textbf{One member's training:} a function that builds the configuration, model,
    optimiser, scheduler and monitor, trains, computes the out-of-bag loss and appends the
    summary row; the training loop is that of \file{heft_train.ipynb}.
  \item \textbf{The loop} over \code{MEMBERS}, skipping finished members.
  \item \textbf{Results on the holdout spectra} (\cref{sec:test}): the ensemble mean
    against POWHEG; spectra with ensemble bands; the spread against $\mhh$; the pull test;
    and, when both ensembles exist, the decomposition (\ref{eq:totvar}) against $\mhh$, and
    the comparison with the BNN run if there is one.
\end{enumerate}

% ===========================================================================
\section{The test: is the spread the size of the error?}
\label{sec:test}

On the 1,000 holdout spectra, at every bin $k$ with POWHEG value $t_k$, POWHEG noise
$\sigma_k$ (as in \file{heft_test.ipynb}), ensemble mean $\bar f_k$ and ensemble standard
deviation $s_k$, form the pull
\begin{equation}
  z_k = \frac{t_k - \bar f_k}{\sqrt{\sigma_k^2 + s_k^2}} .
  \label{eq:pull}
\end{equation}
If the ensemble spread describes the model's error, and the error and the POWHEG noise
are independent, $z_k$ has mean 0 and standard deviation 1, and $|z_k| < 2$ for about
95\% of bins. If the spread is too small, the standard deviation of $z$ is larger than 1.
The same pull, with $s_k$ the posterior standard deviation, is computed for the BNN, so all
are compared on one scale. Note that this test does not depend on the choice of reference
ensemble (\cref{sec:reference}): it asks only whether the stated uncertainty matches the
observed errors on spectra not used in training.

\begin{center}
\small
\begin{tabular}{p{0.38\textwidth}p{0.54\textwidth}}
\toprule
quantity & reported\\
\midrule
accuracy of the ensemble mean & median $\chi^2/N$; fractions above 2 and 5; mean $\chi^2$
  per bin at 255--555 GeV and above; against the single members and the L1 model\\
width & $s_k/\sigma_k$: median, distribution, against $\mhh$\\
calibration & standard deviation of $z$; fraction $|z| < 2$; both against $\mhh$\\
decomposition & optimisation and data terms of (\ref{eq:totvar}) against $\mhh$\\
out-of-bag (bootstrap only) & each member's loss on the rows it did not see, against its
  validation loss\\
\bottomrule
\end{tabular}
\end{center}

\paragraph{What counts as success.}
\begin{enumerate}
  \item The ensemble mean is at least as accurate as a single L1 model, as averaging over
    members usually makes it.
  \item The spread is \emph{informative}: larger where the model is known to be weak (near
    threshold) than elsewhere.
  \item The standard deviation of $z$ is close to 1 above about 555 GeV. Near threshold I
    expect it to exceed 1, because of the shared bias; how much it exceeds 1 measures that
    bias, which is itself a useful number.
\end{enumerate}

% ===========================================================================
\section{Links with the BNN}

\begin{itemize}
  \item \textbf{Starting points for several HMC chains.} The members of the
    \code{'seeds'} ensemble lie in different valleys. Starting four BNN chains from four of
    them makes $\hat R$ across chains a real test for several valleys. This needs a small
    option in \file{heft_train_bnn.ipynb}, \code{START\_FROM}, a list of parameter files,
    one per chain, to be added after this plan.
  \item \textbf{A common yardstick.} The pull (\ref{eq:pull}) puts the BNN band, the two
    ensemble spreads and the POWHEG noise on one scale.
\end{itemize}

% ===========================================================================
\section{Not in this plan}
\begin{itemize}
  \item No change to \file{heft_train.ipynb}, \file{heft_test.ipynb}, \file{hmc.py} or the
    data files.
  \item The spectrum and noise bootstraps of \cref{sec:reference}; the noise bootstrap is
    a natural sensitivity check later, if the results call for one.
  \item The \code{START\_FROM} option of the BNN notebook, and the tempered likelihood,
    are separate steps.
  \item Bootstrap confidence intervals more refined than the spread (percentile, BCa) are
    left out: with $M = 20$ only the standard deviation is well determined.
  \item Nothing is committed to git.
\end{itemize}

% ===========================================================================
\section{Decisions for you}
\label{sec:decisions}

\begin{keyidea}[Decisions before I proceed]
\begin{enumerate}
  \item \textbf{Two ensembles,} \code{'seeds'} and \code{'bootstrap'} (rows), $M = 20$
    each. Agreed, or start with 10?
  \item \textbf{The split of \file{heft_train.ipynb}}, with only the 274,152 training rows
    resampled. Agreed?
  \item \textbf{The training recipe} of \file{heft_train.ipynb}, unchanged (L1, 200,000
    iterations). Agreed?
  \item \textbf{No real-time monitor window} during ensemble training; printed progress
    instead. Agreed?
  \item \textbf{\code{Ensemble} as a new module,} \file{dihiggs/ensemble.py}, leaving
    \file{hmc.py} unchanged for now. Agreed?
  \item \textbf{Folder names:} \file{runs/ENS_seeds/member_01/}, \ldots. Acceptable?
  \item \textbf{This document} replaces the first version at
    \file{docs/ensemble/ensemble_plan.pdf}. Agreed?
\end{enumerate}
\end{keyidea}

% ===========================================================================
\begin{thebibliography}{99}
\bibitem{venn1866} J.~Venn, \emph{The Logic of Chance}, Macmillan (1866).
\bibitem{reichenbach1949} H.~Reichenbach, \emph{The Theory of Probability}, University of
  California Press (1949).
\bibitem{hajek2007} A.~H\'ajek, ``The reference class problem is your problem too'',
  Synthese \textbf{156} (2007) 563.
\bibitem{fisher1956} R.~A.~Fisher, \emph{Statistical Methods and Scientific Inference},
  Oliver and Boyd (1956).
\bibitem{cox1958} D.~R.~Cox, ``Some problems connected with statistical inference'', Ann.\
  Math.\ Statist.\ \textbf{29} (1958) 357.
\bibitem{buehler1959} R.~J.~Buehler, ``Some validity criteria for statistical inferences'',
  Ann.\ Math.\ Statist.\ \textbf{30} (1959) 845.
\bibitem{robinson1975} G.~K.~Robinson, ``Some counterexamples to the theory of confidence
  intervals'', Biometrika \textbf{62} (1975) 155.
\bibitem{birnbaum1962} A.~Birnbaum, ``On the foundations of statistical inference'', J.\
  Amer.\ Statist.\ Assoc.\ \textbf{57} (1962) 269.
\bibitem{berger1988} J.~O.~Berger and R.~L.~Wolpert, \emph{The Likelihood Principle}, 2nd
  ed., IMS Lecture Notes--Monograph Series \textbf{6} (1988).
\bibitem{coxhinkley1974} D.~R.~Cox and D.~V.~Hinkley, \emph{Theoretical Statistics},
  Chapman \& Hall (1974).
\bibitem{cousins1995} R.~D.~Cousins, ``Why isn't every physicist a Bayesian?'', Am.\ J.\
  Phys.\ \textbf{63} (1995) 398.
\bibitem{efron1979} B.~Efron, ``Bootstrap methods: another look at the jackknife'', Ann.\
  Statist.\ \textbf{7} (1979) 1.
\bibitem{efron1993} B.~Efron and R.~J.~Tibshirani, \emph{An Introduction to the
  Bootstrap}, Chapman \& Hall/CRC (1993).
\bibitem{freedman1981} D.~A.~Freedman, ``Bootstrapping regression models'', Ann.\ Statist.\
  \textbf{9} (1981) 1218.
\bibitem{wu1986} C.~F.~J.~Wu, ``Jackknife, bootstrap and other resampling methods in
  regression analysis'', Ann.\ Statist.\ \textbf{14} (1986) 1261.
\bibitem{davison1997} A.~C.~Davison and D.~V.~Hinkley, \emph{Bootstrap Methods and their
  Application}, Cambridge University Press (1997).
\bibitem{gelman1996} A.~Gelman, X.-L.~Meng and H.~Stern, ``Posterior predictive assessment
  of model fitness via realized discrepancies'', Statistica Sinica \textbf{6} (1996) 733.
\bibitem{lakshminarayanan2017} B.~Lakshminarayanan, A.~Pritzel and C.~Blundell, ``Simple
  and scalable predictive uncertainty estimation using deep ensembles'', NeurIPS 2017,
  arXiv:1612.01474.
\bibitem{wilson2020} A.~G.~Wilson and P.~Izmailov, ``Bayesian deep learning and a
  probabilistic perspective of generalization'', NeurIPS 2020, arXiv:2002.08791.
\bibitem{izmailov2021} P.~Izmailov, S.~Vikram, M.~D.~Hoffman and A.~G.~Wilson, ``What are
  Bayesian neural network posteriors really like?'', ICML 2021, arXiv:2104.14421.
\end{thebibliography}

\end{document}
